\subsection{Top-down control and neuromodulation}

In the preceding validation experiments, we configured the four-cell CPG with Si3 neurons as endogenous bursters with elevated intrinsic excitability to facilitate comparison with experimental perturbation protocols.
In the biological circuit, however, Si3 neurons are not intrinsically oscillatory but receive continuous excitatory drive from Si1 command neurons that initiate and sustain swim episodes \cite{sakurai2019command}.
We now examine a five-cell network configuration with Si3 neurons set to quiescent baseline states that require Si1 excitation to become active.

\subsubsection{Si1 gating of network activity}

Si1 neurons function as command-like neurons that provide continuous (non-bursting) excitatory drive to both Si3 neurons, which in turn drive CPG rhythm through the contralateral E/I modules \cite{sakurai2019command}.
This hierarchical organization allows Si1 to gate network state transitions while intrinsic CPG dynamics shape burst structure and left-right alternation.

Gradual Si1 depolarization initiates rhythmogenesis by first recruiting Si3 into high-frequency spiking, then Si2; gradual withdrawal terminates the episode (Fig.~\ref{fig:si1_neuromod}A--B).
The model reproduces the experimentally observed short latency between Si1 activation and Si3 recruitment, demonstrating that time-varying excitatory drive from Si1 regulates network burst period throughout swim episodes.
Abrupt Si1 hyperpolarization rapidly terminates ongoing bursting, and upon release the network resumes rhythmic activity (Fig.~\ref{fig:si1_neuromod}C).
These results establish that Si1 excitatory drive is sufficient to initiate, sustain, and terminate rhythmic bursting.

\subsubsection{Neuromodulatory burst frequency control}

Beyond gating, Si1 neurons regulate burst frequency in an activity-dependent manner: experimental recordings show swim frequency increases with Si1 firing rate, encoding graded frequency information rather than merely on/off state \cite{sakurai2019command}.
Intrinsic neuromodulation within CPG circuits has been demonstrated in \textit{Tritonia diomedea}, a nudibranch in the same clade as \textit{Dendronotus}, where serotonergic swim interneurons dynamically modulate synaptic strengths during motor pattern generation \cite{katz1994dynamic, sakurai2003spike}, providing direct precedent for the modulatory mechanisms we explore here.
We compared two mechanisms by which neuromodulation of the Si3$\to$Si2 excitatory pathway could implement such control (Fig.~\ref{fig:si1_neuromod}C--D).

Presynaptic modulation of the accumulation rate parameter $\alpha$ of the logistic Si3$\to$Si2 excitatory synapse produces smooth, graded frequency changes that scale predictably with modulatory input, enabling fine-tuned frequency adjustments across behavioral contexts.
Postsynaptic modulation of maximal synaptic conductance produces a different frequency-control profile across the same gain scan.

Presynaptic modulation thus provides the smooth, stable frequency control appropriate for a locomotor CPG accommodating developmental scaling and environmental variation.
Postsynaptic modulation's high sensitivity could instead enable rapid behavioral state switching.
These findings motivate experimental tests of where Si1-linked modulation acts within the excitatory pathway.
